\chapter{Pseudocode}

\noindent\textbf{Algorithm 5.2: Regev's Algorithm}
\begin{algorithmic}[1]
\State \textbf{Input}: A composite number $N$, where $N = p\cdot q$, R, T, D
\State \textbf{Output}: primes $p$ and $q$
\State $n \gets \lceil \log_2 N \rceil$
\Comment{Classical preprocessing}
\State $d \gets \lceil \sqrt{n} \rceil$
\State $B[] \gets$ the first $d$ primes, in order
\State $i \gets 0$
\Comment{Quantum}
\While{$i < d$}
    \State Prepare the gaussian. Suppose $amps$ represents the amplitudes
    \State Define two quantum registers, $e[size]$ and $prod$
    \State $e \gets amps[e_i]$
    \State $prod \gets \prod_{i=1}^d B[i]^{2e_i}$ \Comment{Modular Exponentiation}
    \State $u = measure(prod)$ \Comment{This changes the amplitudes in e.}
    \State $e \gets QFT(e)$
    \State $L^*_i = measure(e) / D$
    \State $i \gets i + 1$
\EndWhile
\Comment{The rest of the algorithm is purely classical}
\State Form the matrix M from \textbf{Eq 10.1}
\State $M' = LLL(M)$
\State $M'' = GramSchmidt(M')$
\State $SV$ = The list of all vectors $\mathbf{v}$ in M' where the corresponding M'' vector \textbf{v} has $\lVert \mathbf{v} \rVert <\sqrt{k}2^{k/2}T$  
\For {$SV(j)$ in SV}
    \State $X \gets \prod_{i=1}^{d} B[i]^{SV_j(e_i)}$
    \State $l \gets \gcd(X-1,N) $
    \If  {$l \neq 1$ and $l \neq N$}:
        \State \Return l
    \EndIf
\EndFor
\State \Return 0

\end{algorithmic}
\vspace{2cm}
The gate operations in the following pseudocodes are defined as follows:\\ 
Let $\ket{a}$,$\ket{b}$ be some arbitrary states, $\ket{0}$ and $\ket{1}$ to be the 0 and 1 states respectively, and let $\ket{a_n\dots a_2 a_1}$ be the bitwise representation of state $\ket{a}$. Let $f$ be the function that describes the operation that makes the respective qubits into dirty ancillas for each gate, $0^n$ represent a clean register of size n, and $\ket{g}$ represents a register in some arbitrary state g (essentially an ancilla register that doesn't need to be clean).

\textbf{SWAP} gate: $$\ket{a} \ket{b} \to^{SWAP} \ket{b} \ket{a}$$

\textbf{CMMC} (Controlled Modular Multiplication by a Constant) gate:\\ 
A gate applied sequentially to all $a_i$ from 1 to n performing
$$\ket{a_i} \ket{b} \ket{0}\to \ket{a_i} \ket{(c_i^{a_i} \cdot b \pmod N} \ket{f(c_i,a_i,b,N)}$$
EXPLAIN THIS BETTTER1!!!!!!!!!!!!!!!!!!!!!!!!!!!!!!!!!!!!!!!!!!!!!!!!!!!!!!!!!!!!!!!!!!!!!!!!!!

\textbf{QQMULT} (Quantum-Quantum Modular Multiplication) gate (out of place): $$\ket{a} \ket{b} \ket{t} \ket{0^S} \to^{QQMULT} \ket{a} \ket{b} \ket{(t + a \cdot b) \pmod N} \ket{0^S}$$

\textbf{CPR} (Controlled Phase Rotation) gate: $$\ket{a_k} \ket{b_j} \to^{CPR} \ket{a_k} \ket{exp(2\pi i \frac{a_k}{2^{j - k + 1}} b_j} $$

\textbf{MODPROD} (Modular Product) gate (essentially a form of reversible in-place multiplication)
$$\ket{a}\ket{a^{-1}} \ket{b} \ket{b^{-1} \ket{g} \ket{0^S}} \to^{MODPROD} \ket{a}\ket{a^{-1}} \ket{ab} \ket{(ab)^{-1}} \ket{g} \ket{0^S}$$

\textbf{MSQUARE} (Modular Squaring) gate: $$\ket{a} \ket{a} \to^{MSQUARE} \ket{a} \ket{a^2 \pmod N}$$



\noindent\textbf{Algorithm 5.3: Quantum Fourier Transform}
\begin{algorithmic}[1]
\State \textbf{Input}:Quantum registers consisting of $\ket{e_1,e_2,\dots, e_d}$ in a 
superposition where the $e_i$ are $L$ qubit numbers, $N$.
\State \textbf{Output}: The result of the QFT in the same registers.
\State Let $e_{i,j}$ being the $i_{th}$ qubit of the $j_{th}$ register, from MSB to LSB.
\For{$j$ from 1 to $d$}:
    \State $\ket{e_{i,j}} \to^{H} \ket{e_{i,j}}$
    \For{$i$ from $1$ to $L$}:
            \For {$k$ from $i+1$ to $L$}
                \State $\ket{e_{i,j}}\ket{e_{k,j}} \to^{CPR(e_{k,j},e_{i,j})} (\ket{e_{i,j}})\ket{e_{k,j}}$ 
        \EndFor
    \EndFor
    \For{$i$ from $1$ to $\lfloor L/2 \rfloor$}
        \State $\ket{e_{i,j}}\ket{e_{L-i+1, j}} \to^{SWAP} \ket{e_{i,j}} \ket{e_{L-i+1, j}}$
    \EndFor
\EndFor

\end{algorithmic}

